%!%
% Copyright (c) 2026 mingcheng <mingcheng@outlook.com>
% 
% File: conclusion.tex
% Author: mingcheng <mingcheng@outlook.com>
% File Created: 2026-03-28 09:28:17
% 
% Modified By: mingcheng <mingcheng@outlook.com>
% Last Modified: 2026-03-29 07:17:08
%%

% !TeX root = ../prompt-engineering-guide-for-beginners.tex

\chapter{结束语}

恭喜你读到了这里。

还记得第一章里那个比喻吗？去餐厅点菜，说「来个菜」和说「来一份不辣的番茄炒蛋，少油少盐」，结果天差地别。跟 AI 打交道的道理完全一样。一路读下来你会发现，所谓的提示工程，说到底就是和 AI 好好说话，把话说清楚这件事。角色设定、上下文补充、输出约束、链式推理，这些听起来各有各的名字，但本质上都是在帮你把一句模糊的需求拆成一段 AI 能准确执行的指令。

读完这本小册子只是起点。提示工程是一项需要在实践中打磨的手艺，就像学骑自行车一样，光看教程永远学不会，你得自己骑上去摔几次才能找到平衡感。明天写邮件的时候，整理会议纪要的时候，甚至帮孩子检查作业的时候，不妨打开 AI 试一试。第一次的结果大概率不完美，但没关系，换个角色、加点约束、给几个示例，再跑一遍，你就已经在做真正的提示工程了。每一次这样的迭代，都在把你对 AI 的理解往前推一步。

AI 领域的变化速度远超我们的想象，今天好用的技巧明天可能就有更好的替代方案。但你在这本小册子里学到的那些底层原则，「结构化表达、提供上下文、明确约束」，不会因为模型换代就失效。掌握了这些，新工具出来的时候你也能很快上手。而如果在使用过程中你踩到了有意思的坑，或者发现了一条特别好用的提示词，不妨分享给身边的朋友和同事。教是最好的学，在分享中你自己也会理解得更深。

最后，感谢你花时间读完这本小册子。希望它能成为你和 AI 打交道的一块敲门砖，帮你在工作和生活中真正用好这个时代最强大的工具。
